\documentclass[12pt]{amsart}
% Package Header

\usepackage[english]{babel}
\usepackage{amsfonts}
\usepackage{amsmath}
\usepackage{amssymb}
\usepackage{animate}
\usepackage{array}
\usepackage{arrayjobx}
\usepackage{arydshln}
\usepackage{calculator}
\usepackage{cancel}
\usepackage{enumerate}
\usepackage{mathabx}
\usepackage{mathtools}
\usepackage{multicol}
\usepackage{multimedia}
\usepackage{multirow}
\usepackage{palatino}
\usepackage{polynom}
\usepackage{soul}
\usepackage{tcolorbox}
\usepackage{tikz}
\usepackage{wrapfig}
\usepackage{xcolor}
\usetikzlibrary{arrows,shapes,trees,shapes.geometric,calc}
\usepackage{array}
\usepackage{geometry}
\usepackage{enumitem}
\usepackage{parskip}
\usepackage{booktabs}
\usepackage{graphicx}
\usepackage{mathrsfs}

% \geometry{margin=1in, top=0.75in}
% \topmargin = -1.0in
% \textwidth = 6.6in
% \oddsidemargin = -0.9in
% \evensidemargin = -0.9in

\renewcommand{\headrulewidth}{0pt} % remove lines as well

\numberwithin{equation}{section}

\DeclareGraphicsRule{.tif}{png}{.png}{`convert #1 `dirname #1`/`basename #1 .tif`.png}

\textheight = 10.8 in
\topmargin = -1.0in
\textwidth = 7.7in
\oddsidemargin = -0.7in
\evensidemargin = -0.7in

\begin{document}

\newcommand{\NP}{\newpage\thispagestyle{empty}}
\newcommand{\B}{\textcolor{blue}}
\newcommand{\R}{\textcolor{red}}
\newcommand{\G}{\textcolor{olive}}
\newcommand{\Or}{\textcolor{YellowOrange}}
\newcommand{\DEF}{\underline{Defn}:\ \ \ }
\newcommand{\THM}{\underline{Thm}:\ \ \ }
\newcommand{\PROP}{\underline{Properties}:\ \ \ }
\newcommand{\NOTE}{\underline{Notes}:\ \ \ }
\newcommand{\EX}{\underline{Ex}:\ \ \ }
\newcommand{\SOL}{\underline{Soln}:\ \ \ }
\newcommand{\UL}{\underline}
\newcommand{\RE}{\mathbf{R}}
\newcommand{\V}{\textbf{ v }}
\newcommand{\VV}{[v_1,v_2]}
\newcommand{\VVV}{[v_1,v_2,v_3]}
\newcommand{\W}{\textbf{ w }}
\newcommand{\WW}{[w_1,w_2]}
\newcommand{\WWW}{[w_1,w_2,w_3]}
\newcommand{\VA}{\textbf{ a }}
%\newcommand{\AA}{[a_1,a_2]}
\newcommand{\AAA}{[a_1,a_2,a_3]}
\newcommand{\VB}{\textbf{ b }}
\newcommand{\BB}{[b_1,b_2]}
\newcommand{\BBB}{[b_1,b_2,b_3]}
\newcommand{\VC}{\textbf{ c }}
\newcommand{\CC}{[c_1,c_2]}
\newcommand{\CCC}{[c_1,c_2,c_3]}
\newcommand{\VI}{\textbf{ i }}
\newcommand{\VJ}{\textbf{ j }}
\newcommand{\VK}{\textbf{ k }}
\newcommand{\VF}{\textbf{ F }}
\newcommand{\VG}{\textbf{ G }}
\newcommand{\VO}{{\vec{0}}}
\newcommand{\OBUL}{\quad $\circ$ \ \ \ }
\newcommand{\NEWLINE}{\\ \hspace{0.63cm}}
\newcommand{\CAN}{\cancel}
\newcommand{\GF}{{\nabla f}}
\newcommand{\PFPX}{{\dfrac{\partial f}{\partial x}}}
\newcommand{\PFPY}{{\dfrac{\partial f}{\partial y}}}
\newcommand{\PFPZ}{{\dfrac{\partial f}{\partial z}}}
\newcommand{\DFDX}{{\dfrac{df}{dx}}}
\newcommand{\GG}{{\nabla g}}
\newcommand{\PGPX}{{\dfrac{\partial g}{\partial x}}}
\newcommand{\PGPY}{{\dfrac{\partial g}{\partial y}}}
\newcommand{\PGPZ}{{\dfrac{\partial g}{\partial z}}}
\newcommand{\DGDX}{{\dfrac{dg}{dx}}}
\newcommand{\PPX}{{\dfrac{\partial}{\partial x}}}
\newcommand{\PPY}{{\dfrac{\partial}{\partial y}}}
\newcommand{\PPZ}{{\dfrac{\partial}{\partial z}}}
\newcommand{\PP}{\partial}
\newcommand{\DELF}{\nabla \cdot \textbf{ F }}
\newcommand{\CURLF}{\nabla \times \textbf{ F }}
\newcommand{\DEL}{\nabla \cdot }
\newcommand{\CURL}{\nabla \times }
\newcommand{\RT}{\textbf{ r }(t)}
\newcommand{\RPT}{\textbf{ r }'(t)}
\newcommand{\GINT}{\iint \limits_{R} (\text{curl} \ \textbf{ F }) \cdot \textbf{ k } \ dA}
\newcommand{\ON}{\onslide}
\newcommand{\OB}{\overbrace}
\newcommand{\OBT}{\overbracket}
\newcommand{\OS}{\overset}
\newcommand{\UB}{\underbrace}
\newcommand{\UBT}{\underbracket}
\newcommand{\US}{\underset}
\newcommand{\D}{\displaystyle}
\newcommand{\MCol}{\multicolumn}
\newcommand{\RRE}{\mathbf{R}^2}
\newcommand{\RRRE}{\mathbf{R}^3}
\newcommand{\DXDT}{{\dfrac{dx}{dt}}}
\newcommand{\DYDT}{{\dfrac{dy}{dt}}}
\newcommand{\DZDT}{{\dfrac{dz}{dt}}}
\newcommand{\FRT}{\textbf{ F }(\textbf{ r }(t))}
\newcommand{\WORK}{ {\int \limits_{C} \textbf{ F } \cdot d\textbf{ r }} }
\newcommand{\WORKAB}{ {\int \limits_{A}^{B}\textbf{ F } \cdot d\textbf{ r }} }
\newcommand{\WORKT}{ {\int \limits_{a}^{b}\textbf{ F }(\textbf{ r }(t)) \cdot \textbf{ r }'(t) \ dt} }
\newcommand{\FLUX}{ {\iint \limits_{S}\textbf{ F } \cdot d\textbf{ S }} }
\newcommand{\FLUXX}{ {\iint \limits_{S}\textbf{ F } \cdot \textbf{ n } \ dS} }
\newcommand{\FLUXXX}{ {\iint \limits_{R}\textbf{ F }(\textbf{ r }(u,v)) \cdot \textbf{ N } \ du\ dv} }
\newcommand{\OWORK}{ {\oint \limits_{C}\textbf{ F } \cdot d\textbf{ r }} }
\newcommand{\GREENT}{ {\oint \limits_{C}\textbf{ F }\cdot \textbf{ T } \ ds} }
\newcommand{\GREENN}{ {\oint \limits_{C}\textbf{ F }\cdot \textbf{ N } \ ds} }
\newcommand{\DIVTHM}{\iiint \limits_{T} \text{div} \ \textbf{ F }  \ dV}
\newcommand{\GREENCURL}{\iint \limits_{R} (\text{curl} \ \textbf{ F }) \cdot \textbf{ k } \ dA}
\newcommand{\STOKECURL}{\iint \limits_{S} (\text{curl} \ \textbf{ F }) \cdot \textbf{ n } \ dS}
\newcommand{\GREENDIV}{\iint \limits_{R} \text{div} \ \textbf{ F } \ dA}
\newcommand{\GREENSA}{\iint \limits_{R} \left(\dfrac{\partial N}{\partial x} - \dfrac{\partial M}{\partial y}\right) \ dA}
\newcommand{\GREENSL}{\oint \limits_{C} (M\ dx + N\ dy )}
\newcommand{\FDR}{ {\textbf{ F } \cdot d\textbf{ r }} }
\newcommand{\FDRT}{ {\textbf{ F }(\textbf{ r }(t)) \cdot \textbf{ r }'(t)} }
\newcommand{\BCOLS}{\begin{columns}}
\newcommand{\BCOL}{\begin{column}}
\newcommand{\ECOLS}{\end{columns}}
\newcommand{\ECOL}{\end{column}}
\newcommand{\BLINE}{\tikz[baseline]}
\newcommand{\NODERRB}{\node[fill=red!10,draw,rectangle,anchor=base]}
\newcommand{\NODEGRB}{\node[fill=green!10,draw,rectangle,anchor=base]}
\newcommand{\NODEBRB}{\node[fill=blue!10,draw,rectangle,anchor=base]}
\newcommand{\NODECRB}{\node[fill=red!0,draw,rectangle,anchor=base]}
\newcommand{\NODERR}{\node[fill=red!10,rectangle,anchor=base]}
\newcommand{\NODEGR}{\node[fill=green!10,rectangle,anchor=base]}
\newcommand{\NODEBR}{\node[fill=blue!10,rectangle,anchor=base]}
\newcommand{\NODECR}{\node[fill=red!0,rectangle,anchor=base]}
\newcommand{\NODERE}{\node[fill=red!10,draw,ellipse,anchor=base]}
\newcommand{\NODEGE}{\node[fill=green!10,draw,ellipse,anchor=base]}
\newcommand{\NODEBE}{\node[fill=blue!10,draw,ellipse,anchor=base]}
\newcommand{\NODECE}{\node[fill=red!0,draw,ellipse,anchor=base]}
\newcommand{\NODERC}{\node[fill=red!10,draw,circle,anchor=base]}
\newcommand{\NODEGC}{\node[fill=green!10,draw,circle,anchor=base]}
\newcommand{\NODEBC}{\node[fill=blue!10,draw,circle,anchor=base]}
\newcommand{\NODECC}{\node[fill=blue!0,draw,circle,anchor=base]}
\newcommand{\CNODE}{\tikz[na]\node [coordinate] }
\newcommand{\BENDL}{edge [bend left] }
\newcommand{\BENDR}{edge [bend right] }
\newcommand{\ARROWR}{\path [red,->]}
\newcommand{\ARROWG}{\path [green,->]}
\newcommand{\ARROWB}{\path [blue,->]}
\newcommand{\ARROWSR}{\path [red,<->]}
\newcommand{\ARROWSG}{\path [green,<->]}
\newcommand{\ARROWSB}{\path [blue,<->]}
\newcommand{\LINER}{\path [red,-]}
\newcommand{\LINEG}{\path [green,--]}
\newcommand{\LINEB}{\path [blue,--]}
\newcommand{\DE}{east}
\newcommand{\DW}{west}
\newcommand{\DN}{north}
\newcommand{\DS}{south}

\newcommand{\TT}[1]{\texttt{#1}}

\tikzstyle{every picture}+=[remember picture]
\tikzstyle{na} = [baseline=-1.5ex]

%%-----------------------------------------------------------
%% Bliss Definitions
\newcommand{\To}{\Rightarrow}
\newcommand{\del}{\nabla}
\newcommand{\C}{\mathcal{C}}
\newcommand{\N}{\mathcal{N}}
\newcommand{\A}{\mathcal{A}}
\renewcommand{\O}{\mathcal{O}}
% \renewcommand{\N}{\mathbb{N}}
\newcommand{\ben}{\begin{enumerate}}
\newcommand{\een}{\end{enumerate}}
\newcommand{\eps}{\varepsilon}
\renewcommand{\epsilon}{\varepsilon}
\newcommand{\p}{\partial}
\newcommand{\<}{\langle}
\renewcommand{\>}{\rangle}
\renewcommand{\subset}{\subseteq}
\newcommand{\cont}{\Rightarrow \Leftarrow}
\newcommand{\back}{\backslash}
\newcommand{\abs}[1]{|{#1}|}
\newcommand{\ip}[1]{\langle{#1}\rangle}
\newcommand{\m}[1]{{\mathbf{#1}}}
\newcommand{\sts}{\m{S}^\textrm{T}\m{S}}
\newcommand{\rank}{\textrm{rank}}
\newcommand{\X}{\mathcal{X}}
\newcommand{\Y}{\mathcal{Y}}
\def\ds{\displaystyle}
\def\d{\partial}
\renewcommand{\arraystretch}{1.5}
%%-----------------------------------------------------------


\newcommand\CEN[1]{\multicolumn{1}{c}{ #1 }}

\newcommand\MATL[3]
{{%
\arraycolsep=#1\arraycolsep\ensuremath{{\setlength{\extrarowheight}{#2cm}\begin{bmatrix*}[l]#3\end{bmatrix*}}}
}}
\newcommand\MATC[3]
{{%
\arraycolsep=#1\arraycolsep\ensuremath{{\setlength{\extrarowheight}{#2cm}\begin{bmatrix*}[c]#3\end{bmatrix*}}}
}}
\newcommand\MATR[3]
{{%
\arraycolsep=#1\arraycolsep\ensuremath{{\setlength{\extrarowheight}{#2cm}\begin{bmatrix*}[r]#3\end{bmatrix*}}}
}}

\newcommand\MATDL[3]
{{%
\arraycolsep=#1\arraycolsep\ensuremath{{\setlength{\extrarowheight}{#2cm}\begin{vmatrix*}[l]#3\end{vmatrix*}}}
}}
\newcommand\MATDC[3]
{{%
\arraycolsep=#1\arraycolsep\ensuremath{{\setlength{\extrarowheight}{#2cm}\begin{vmatrix*}[c]#3\end{vmatrix*}}}
}}
\newcommand\MATDR[3]
{{%
\arraycolsep=#1\arraycolsep\ensuremath{{\setlength{\extrarowheight}{#2cm}\begin{vmatrix*}[r]#3\end{vmatrix*}}}
}}

\newcommand\TABwB[3]
{{
\setlength{\extrarowheight}{#1 cm}
\begin{tabular}{#2}
	\hline
	#3
	\hline
\end{tabular}
}}

\newcommand\TABwoB[3]
{{
\setlength{\extrarowheight}{#1cm}
\begin{tabular}{#2}
	#3
\end{tabular}
}}

\newcommand\includegraphicsDD[2]
{%
\includegraphics[scale=#1]{#2}
}

\newcommand\includegraphicsDDD[2]
{%
\includemedia[scale=#1,3Dmenu,activate=onclick,deactivate=onclick,
3Droll=0.,
3Dortho=0.0046491301618516445,
3Dc2c=0.8617081642150879 -0.29301872849464417 0.4142450988292694,
3Dcoo=-13.165555000305176 3.489872455596924 52.97535705566406,
3Droo=149.99999574342704,
3Dlights=Headlamp,add3Djscript=asylabels.js]{\includegraphics{#2}}
{#2.prc}
}

\newcommand\MATS[2]
{{%
  \arraycolsep=#1\arraycolsep\ensuremath{\begin{bmatrix*}[r]#2\end{bmatrix*}}
}}

\newcommand{\OKFACE}{

\begin{tikzpicture}
\draw[thick] (0cm,0cm) circle(.2cm);
\draw[thick] (-0.1,-0.09) -- (0.1,-0.09);
\draw [thick, fill=black] (0,0) circle (0.01);
\draw [rotate=90, fill=white] (0.06,0.06) circle (0.03);
\draw [rotate=90, fill=white] (0.06,-0.06) circle (0.03);
\end{tikzpicture}
}

\newcommand{\HAPPYFACE}{

\begin{tikzpicture}
\draw[thick] (0cm,0cm) circle(.2cm);
\draw[thick] plot [smooth,tension=1.5] coordinates{(-0.1,-0.08) (0,-0.13) (0.1,-0.08)};
\draw [thick, fill=black] (0,0) circle (0.01);
\draw [rotate=90, fill=white] (0.06,0.06) circle (0.04);
\draw [rotate=90, fill=white] (0.06,-0.06) circle (0.04);
\end{tikzpicture}
}

\newcommand{\SADFACE}{

\begin{tikzpicture}
\draw[thick] (0cm,0cm) circle(.2cm);
\draw[thick] plot [smooth,tension=1.5] coordinates{(-0.1,-0.12) (0,-0.07) (0.1,-0.12)};
\draw [thick, fill=black] (0,0) circle (0.01);
\draw [rotate=90, fill=white] (0.06,0.06) circle (0.02);
\draw [rotate=90, fill=white] (0.06,-0.06) circle (0.02);
\end{tikzpicture}
}

\newcommand{\CounterReset}{
\setcounter{CritCounter}{1}
\setcounter{TestCounter}{1}
\setcounter{ColorCounter}{1}
\setcounter{Index}{1}
}

\newcommand\CritLineInitialize[1]{
\ifdef{\SymFun}
{
	\setcounter{CritLineSlide}{#1}
}
{
	\newarray\TestPts
	\newarray\TestVals
	\newarray\TestEVal
	\newarray\fIncDec
	\newarray\testsigns
	\newarray\inequality
	\setcounter{CritLineSlide}{#1}
	\newcommand\SymFun{}
	\newcommand\IncDecStr{}
	\newcommand\TestSignStr{}
	\newcommand\InEqStr{}
	\newcommand\SymFunD{}
	\newcommand\SymFunDLt{}
	\newcommand\SymFunN{}
	\newcommand\SymFunNLt{}
	\newcommand{\CritLineTable}{}
}
}

\newcommand\USubInitialize{
	\newcommand{\Fun}{}
	\newcommand{\FunN}{}
	\newcommand{\FunD}{}
	\newcommand{\FunU}{}
	\newcommand{\UIntegrand}{}
	\newcommand{\AFunU}{}
	\newcommand{\DX}{}
}

\newcommand*\circled[1]{\tikz[baseline=(char.base)]{
            \node[shape=circle,draw,inner sep=1pt,color=blue] (char) {#1};}}

\newcommand*\ellipsed[1]{\tikz[baseline=(char.base)]{
            \node[shape=ellipse,draw,inner sep=1pt,color=blue] (char) {#1};}}

\newcommand*\rected[1]{\tikz[baseline=(char.base)]{
            \node[shape=rectangle,draw,inner sep=1pt,color=blue] (char) {#1};}}

\newcommand*\starred[1]{\tikz[baseline=(char.base)]{
            \node[shape=star,star points=10,draw,inner sep=1pt,color=blue] (char) {#1};}}

\newcommand\CHECKBOX[1]{\makebox[0pt][l]{$\square$}\raisebox{.15ex}{\hspace{0.1em}\ON#1{\B{$\checkmark$}}}}

\makeatletter
\newcommand{\@RemovePositiveSignIfAny}[1]{\IfBeginWith{#1}{+}{\StrGobbleLeft{#1}{1}}{#1}}%
\newcommand{\@PlusIfNonZero}[1]{\IfStrEq{#1}{}{}{+}}%
\newlength{\@MySystemWidth}%
\newcommand*{\@ChooseSignWithAlignment}[2]{%
    \IfBeginWith{#2}{-}%
    {{}-\makebox[\@MySystemWidth][#1]{\ensuremath{\StrGobbleLeft{#2}{1}}}}%
    {{}\@PlusIfNonZero{#2}\makebox[\@MySystemWidth][#1]{\ensuremath{\@RemovePositiveSignIfAny{#2}}}}%
}%
\newcommand*{\@ChooseSignWithRightAlignment}[1]{\@ChooseSignWithAlignment{r}{#1}}%
\newcommand*{\@ChooseSignWithLeftAlignment}[1]{\@ChooseSignWithAlignment{l}{#1}}%

\newenvironment{MySystem}[3][1.2em]{% [<width>] {num of variables} {alignment}
    \newcolumntype{R}{>{\collectcell\@ChooseSignWithRightAlignment}{r}<{\endcollectcell}}
    \newcolumntype{L}{>{\collectcell\@ChooseSignWithLeftAlignment}{l}<{\endcollectcell}}
    \setlength{\@MySystemWidth}{#1}% Store the width to use to typeset the terms
    \IfStrEqCase{#3}{%
        {R}{\begin{array}{r@{} *{\numexpr#2-1\relax}{#3@{}} @{{}={}} r}}%
        {L}{\begin{array}{l@{} *{\numexpr#2-1\relax}{#3@{}} @{{}={}} l}}%
    }[\parbox{\linewidth}{\color{red}MySystem Error: Only L and R alignment types are supported. (using R)}\begin{array}{r *{\numexpr#2-1\relax}{#3} @{{}={}} r}]%
}{%
    \end{array}%
}
\newenvironment{MySystemNoEqual}[3][1.2em]{% [<width>] {num of variables} {alignment}
    \newcolumntype{R}{>{\collectcell\@ChooseSignWithRightAlignment}{r}<{\endcollectcell}}
    \newcolumntype{L}{>{\collectcell\@ChooseSignWithLeftAlignment}{l}<{\endcollectcell}}
    \setlength{\@MySystemWidth}{#1}% Store the width to use to typeset the terms
    \IfStrEqCase{#3}{%
        {R}{\begin{array}{r@{} *{\numexpr#2\relax}{#3@{}} @{{}{}} r}}%
        {L}{\begin{array}{l@{} *{\numexpr#2\relax}{#3@{}} @{{}{}} l}}%
    }[\parbox{\linewidth}{\color{red}MySystem Error: Only L and R alignment types are supported. (using R)}\begin{array}{r *{\numexpr#2\relax}{#3} @{{}{}} r}]%
}{%
    \end{array}%
}

\makeatother


$\quad$

\vspace{0.6cm}

\begin{center}
\LARGE {MA 301} \\

\Large{Final} \\

\large{Fall - 2014} \\

\vspace{0.5cm}

{\setlength{\extrarowheight}{0.2cm}
\begin{tabular}{|c|c|c|}
\hline
Problem & points & out of  \\
\hline
\hline
1 & & 6 \\
\hline
2 & & 6 \\
\hline
3 & & 6 \\
\hline
4 & & 6 \\
\hline
5 & & 6 \\
\hline
6 & & 6 \\
\hline
7 & & 16 \\
\hline
8 & & 14 \\
\hline
\hline
Score & & 66\\
\hline
\end{tabular}}
\vspace{0.5cm}

%--------TABLE---------
	{\setlength{\extrarowheight}{0.2cm}
	\begin{tabular}{rcc}
		Name: & \underline{$\qquad \qquad \qquad \qquad \quad$}& \\
		&&$\hspace{1.4cm}$\\
		HR: & \underline{$\qquad \qquad \qquad \qquad \quad$}& \\
		&&\\
		Signature: & \underline{$\qquad \qquad \qquad \qquad \quad$}& \\
	\end{tabular}}
%--------------------

\vspace{1.0cm}

	{\setlength{\extrarowheight}{0.2cm}
	\begin{tabular}{| p{12cm} |}
		\hline
		\multicolumn{1}{|c|}{{\bf READ BEFORE YOU START}} \\
		\hline
		 * Show all work.\\
		 * Change your calculator mode to RADIANS.\\
		 * Sloppy solutions may be penalized.\\
		 * Use calculators ONLY for computations.\\
		 * You must write up every problem in a clear and \\
		 $\,\,$ organized way.\\
		 * If you write solutions on a page different than the \\
		 $\,\,$ question, then clearly state where it is located.\\
		\hline
	\end{tabular}}

\end{center}

\vspace{1.0cm}

\thispagestyle{empty}
\newpage
\thispagestyle{empty}

\noindent
{\bf Problem 1:} Let $\D{{\bf F} = [16xy + 3yx^2, x^3 + 2y]}$ and $C$ is the boundary of the region between:
$x=3$, $y=0$, \& $y=x$.

\vspace{0.2cm}

Compute \, $\D{\OWORK}$

\thispagestyle{empty}
\newpage
\thispagestyle{empty}

\noindent
{\bf Problem 2:} Let ${\bf F} = [-y,2z,0]$ and $C$ be the boundary of side of the half cylinder $y^2+z^2=4$,

 with $z \ge 0$,  \, $0 \le x \le 3$ and oriented outward (see below). 
 
\vspace{0.2cm}

{\setlength{\extrarowheight}{-0.1cm}
\begin{tabular}{p{12.0cm}p{4.0cm}}
	&
	\multirow{2}{*}{\includegraphics[scale = 0.25]{../../../Figures/10_8_Fig3.png}} \\
	Use Stokes' thm to compute \, $\D \WORK$ & \\
	&\\
\end{tabular}}

\thispagestyle{empty}
\newpage
\thispagestyle{empty}

\noindent
{\bf Problem 3:} Given $\D{{\bf F} = [e^x\cos y + yz,\, xz - e^x\sin y,\, xy + z]}$

\vspace{0.2cm}

(a) Show that $\text{\bf F}$ is conservative

%	\begin{tabular}{p{10cm}|p{10cm}}
%		\multicolumn{1}{c|}{method 1}&\multicolumn{1}{c}{method 2}\\
%		&\\&\\&\\&\\&\\&\\&\\&\\&\\&\\&\\&\\&\\&\\&\\&\\&\\&\\&\\&\\&\\&\\&\\&\\&\\&\\&\\&\\&\\&\\&\\&\\&\\&\\&\\&\\&\\
%	\end{tabular}

\vspace{15.4cm}

(b) Evaluate $\displaystyle{\int \limits_{C} \text{\bf F} \cdot  d\text{\bf r}}$ when $C$ is a curve from $(0,0,2)$ to $\D{\left(0,\frac{3\pi}{2},-6\right)}$

\thispagestyle{empty}
\newpage
\thispagestyle{empty}

\noindent
{\bf Problem 4:} Let ${\bf F} = [x^3-y^3,y^3-z^3,z^3-x^3]$ and let 

%-------NUM LIST---------
\begin{enumerate}
	\item[$i.$] \,\, $S_S$ be the upper-half sphere: $x^2 + y^2 + z^2 = 25$, \, $z \ge 0$
	\item[$ii.$] \,\, $S_C$ the region bounded by circle $x^2+y^2=5$ in the $x,y$-plane.
	\item[$iii.$] \,\, $S = S_S \cup S_C$
\end{enumerate}
%------------------------

\noindent
Assuming outward orientation, Compute \,\, $\D \iint \limits_{S} {\bf F} \cdot {\bf n} \, dS$

\vspace{0.2cm}

\thispagestyle{empty}
\newpage
\thispagestyle{empty}

\noindent
{\bf Problem 5:} Let \,\,$z_1 = -2+2i$,\,\, $z_2 = -i$,\,\, $z_3 = 1-i$. \quad Answer the following in the ({\em specified form}):

\vspace{0.2cm}

(a) ({\em Standard Form}) $z_1z_2 = $

\vspace{2.8cm}

(b) ({\em Standard Form}) $z_3-z_1z_2 = $

\vspace{2.8cm}

(c) Convert $z_3$ to Euler's Polar form:

\vspace{3.8cm}

(d) ({\em Euler's Polar Form})  $(z_3)^{76} = $

\vspace{3.8cm}

(e) ({\em Euler's Polar Form}) Solve for $x$: \quad $\D{\frac{1}{\sqrt{2}}x^3 - z_3 = 0}$ \quad (explicitly write out all solutions)

\thispagestyle{empty}
\newpage
\thispagestyle{empty}

\noindent
{\bf Problem 6:} Let \quad
$
f(x) = 
\left\{
\begin{array}{lrc}
	 \pi & & -\frac{\pi}{2} < x < 0\\
	-\pi & & 0  < x < \frac{\pi}{2}\\
\end{array}
\right.
$  \quad and \quad $f(x) = f(x+\pi)$

\vspace{0.1cm}

Find the Fourier Series of $f(x)$ 

\thispagestyle{empty}
\newpage
\thispagestyle{empty}

%{\bf Problem 8:} Let \quad
%$
%f(x) = 
%\left\{
%\begin{array}{lrc}
%	0 & & -1 < x < 0\\
%	e^{-\pi x} & & 0 < x < 1\\
%\end{array}
%\right.
%$  \quad and \quad $f(x) = f(x+2)$
%
%\vspace{0.1cm}
%
%Find the Fourier Series of $f(x)$
%
%\thispagestyle{empty}
%\newpage
\thispagestyle{empty}

\noindent
{\bf Problem 7:} Answer the following true/false and short-answer questions:

\vspace{0.1cm}

{{\setlength{\extrarowheight}{-0.4cm}
\begin{tabular}{|cp{16.4cm}|}
	\hline
	&\\
	\, (a) \, \fbox{\it T} \fbox{\it F} 
	& If ${\bf F}$ is path-independent, then there is a potential function for ${\bf F}$\\
	&\\
	\hline
	&\\
	\, (b) \, \fbox{\it T} \fbox{\it F} 
	& If ${\bf F}$ is path-dependent, then $\WORK \ne 0$ for some closed path $C$.\\
	&\\
	\hline
	&\\
	\, (c) \, \fbox{\it T} \fbox{\it F} 
	& If ${\bf F}$ is path-independent, then $\WORK = 0$ for some closed path $C$.\\
	&\\
	\hline
	&\\
	\, (d) \, \fbox{\it T} \fbox{\it F} 
	& Let $S$ be a sphere of radius 1 centered at the origin and oriented outward.  If $\FLUX = 0$, then ${\bf F} = [0,0,0]$\\
	&\\
	\hline
\end{tabular}}

{{\setlength{\extrarowheight}{-0.4cm}
\begin{tabular}{|cp{14.54cm}|}
	&\\
	\, (g) \, {\it Fill in the blank: } & 
	$\sqrt{-5}\cdot\sqrt{-5} = \UL{\hspace{2.0cm}}$\\
	&\\
	\hline
	&\\
	\, (h) \, {\it Fill in the blank: } & 
	Suppose $f(x) = f(x+6)$ for all $x$, then we say $f$ has a $\text{period of } \UL{\hspace{2.0cm}}$\\
	&\\
	\hline
\end{tabular}}

{{\setlength{\extrarowheight}{-0.4cm}
\begin{tabular}{|cp{17.74cm}|}
	&\\
	\, (e)& 
	Let $S$ be the surface of a cube centered at the origin with side length $h$ and outward orientation.  If ${\bf F} = x\VI + xz\VJ - 3z\VK$ and the flux of ${\bf F}$ through $S$ is $-16$.  Find $h$.\\
	&\\&\\&\\&\\&\\&\\&\\&\\&\\&\\&\\&\\&\\&\\&\\
	\hline
	&\\
	\, (f)& 
	If ${\bf F} = [xz,3,y^3], \, \iint \limits_{S} {\bf F} \cdot {\bf n} \, dS = 12\pi$ and $S$ is a disk of radius $r$ in the $xz$-plane. Find $r$.\\
	&\\&\\&\\&\\&\\&\\&\\&\\&\\&\\&\\&\\&\\&\\&\\
	\hline
\end{tabular}}

{{\setlength{\extrarowheight}{-0.4cm}
\begin{tabular}{| p{18.82cm} |}
	\\
	\begin{tabular}{p{13.8cm}c}
		&
		\multirow{2}{*}{\includegraphics[scale = 0.21]{../../../Figures/10_8_Fig8.png}} \\
		(i) Let $\D{{\bf F}(x,y,z) = [zy-3y,xz-z,xy+1]}$. \, Find \, $\D \int \limits_{C} {\bf F} \cdot d{\bf r} $
		&\\
	\end{tabular} \\
	\\\\\\\\\\\\\\\\\\\\\\\\\\\\\\\\
	\hline
	\begin{tabular}{p{12.4cm}c}
		&\\
		& \multirow{1}{*}{\includegraphics[scale = 0.23]{../../../Figures/10_8_Fig6.png}}\\
		(j) Let $\D{{\bf F} = \nabla\left(\frac{y-1}{y+z+1}\right)}$ and let\, $C$ and $D$ be shown in the figure. &\\
		\\
		$\D{\int \limits_{C} {\bf F} \cdot d{\bf r} = }$ \\
		&\\
		\\\\\\\\
		$\D{\int \limits_{D} {\bf F} \cdot d{\bf r} = }$ &\\
		\\\\\\\\
	\end{tabular}
	\\\\
	\hline
	\\
	(k) Let $\text{\bf F} = [6x,ye^z,-e^z]$ and $S$ be the surface (oriented outward) of the box $|x|\le 1$, $|y|\le 2$ , $0\le z \le 1$.
	\\\\
	Find \, $\D{\iint \limits_{S} \text{\bf F}\cdot \text{\bf n} \, dA =}$
	\\\\\\\\\\\\\\\\\\\\\\\\\\\\
	\hline
\end{tabular}}

\thispagestyle{empty}
\newpage
\thispagestyle{empty}

\noindent
{\bf Problem 8:} Consider the vector field $\D{{\bf F}(x,y,z) = \left[x,y+3,4+2z\right]}$ and let $S_P$ be part of the paraboloid $z=x^2+y^2$ (front surface), $S_0$ (back surface), $S_T$ (top surface), and $C$ (curve around top) be as shown in the figure.  Let $S = S_P \cup S_0 \cup S_T$.  Finally, let $T$ be the solid bounded by $S$.

{\setlength{\extrarowheight}{0.4cm}
\begin{tabular}{rp{11cm}p{4cm}}
	$Notes$: & & \multirow{4}{*}{\includegraphics[scale = 0.44]{../../../Figures/FinalFig_02.png}}\\
	1. & Assume outward orientation &\\
	2. & $S$ is closed &\\
	3. & $\D \text{Area}(S_0) = 2\int \limits_{0}^{2} \left(4-x^2\right) \, dx = \frac{32}{3}$ &\\
	4. & $\D \text{Area}(S_T) = \frac{1}{2}\left(\pi (2)^2\right) = 2\pi$ &\\
	5. & $\D \text{Volume}(T) = 4 \pi$ &\\
\end{tabular}}

\vspace{0.2cm}

{\setlength{\extrarowheight}{0.1cm}
\begin{tabular}{p{0.35cm}p{20.6cm}}
	&\\
	(a) & Is $\D {\bf F}$ conservative? Justify your answer.\\
	&\\&\\&\\&\\&\\&\\&\\&\\
	(b) & Which major theorem guarantees that $\D \iint \limits_{S_P} \text{curl}({\bf F}) \cdot {\bf n} \, d{S} = \iint \limits_{S_T} \text{curl}({\bf F}) \cdot {\bf n} \, d{S}$. \\
	&\\&\\&\\&\\
	(c) & $\D \iint \limits_{S_T} \text{curl}({\bf F}) \cdot {\bf n} \, d{S}$ \\
	&\\&\\&\\&\\&\\&\\&\\&\\&\\&\\&\\
	(d) & Without doing any more work, give $\D \int \limits_{C} {\bf F} \cdot  d{\bf r} = $ \\
	&\\
\end{tabular}}

\thispagestyle{empty}
\newpage
\thispagestyle{empty}

{\setlength{\extrarowheight}{0.1cm}
\begin{tabular}{p{0.35cm}p{20.6cm}}
	(e) & $\D \iint \limits_{S_0} {\bf F} \cdot {\bf n} \, d{S}$ \\
	&\\&\\&\\&\\&\\&\\&\\
	(f) & $\D \iint \limits_{S_P} {\bf F} \cdot {\bf n} \, d{S}$\\
	&\\&\\&\\&\\&\\&\\&\\&\\&\\&\\&\\&\\&\\&\\&\\&\\&\\&\\&\\&\\&\\&\\&\\
	(g) & $\D \iint \limits_{S} {\bf F} \cdot {\bf n} \, d{S}$ \\
	&\\&\\&\\&\\&\\&\\&\\
	(h) & Using (e)-(g), deduce the value of $\D \iint \limits_{S_T} {\bf F} \cdot {\bf n} \, d{S}$ \\
	&\\
\end{tabular}}


\end{document}